% \section{Important Final Considerations}
% \subsection{Designing evaluations based on actual user behavior - should shift from focusing mainly on simulated, one-to-two-turn interactions to longer multi-turn interactions, and studies should cover the variety of ways in which users use LLMs}

% \subsection{Equity: sycophancy may harm some communities, minoritized groups, or populations underrepresented in the training data. For instance, homogenization of writing styles to a more Western style or younger users whose brains are still developing. This should be carefully studied}

% \subsection{Importance of studying downstream effects at both a macro and a micro level}
\section{Conclusion}

In this work, we proposed a unified construct of AI sycophancy, where sycophantic behaviors in AI systems are characterized by displacement of one or more context-relevant norms associated with observed approval-seeking behavior in model generations. This definitions allows researchers to empirically determine which behaviors should be considered sycophantic rather than attempt to redefine the concept, reconcile divergent definitions, or anthropomorphize AI systems. 

